\exercisesection{}

\begin{enumerate}
\def\labelenumi{\arabic{enumi})}
\tightlist
\item
  设 \(\mathbf{T}\) 为一个拓扑空间，\(\mu\) 为 \(\mathbf{T}\) 上的一个正测度。用 \(\mathfrak{S}_{\mu}\) 表示 \(\mathbf{T}\) 中边界为 \(\mu\) -可忽略的子集所成的集合。
\end{enumerate}

\begin{enumerate}
\def\labelenumi{\alph{enumi})}
\item
  证明 \(\mathfrak{S}_{\mu}\) 是一个集环。
\item
  若 T 是完全正则的，则 T 的每个点都有一个包含于 \(S_{\mu}\) 的基本邻域系（注意：若 f 是 T 上的连续函数且在一个 \(\mu\) -可积集合之外为零，则使 \(f(r)\) 不是 \(\mu\) -可忽略的实数 r 所成的集合是可数的）。
\end{enumerate}

\begin{enumerate}
\def\labelenumi{\arabic{enumi})}
\setcounter{enumi}{1}
\tightlist
\item
  设 \(\mathbf{T}\) 是一个完全正则空间，\(\mathfrak{F}\) 是 \(\mathcal{M}_{+}^{b}(\mathrm{T})\) 上的一个滤子，它紧收敛到一个有界测度 \(\mu\)。称 Borel 子集 \(\mathbf{A}\)（\(\mathbf{T}\) 的）为（对 \(\mathfrak{F}\) 的）收敛集，如果 \(\lim_{\lambda,\mathfrak{F}} \lambda_{\mathbf{A}} = \mu_{\mathbf{A}}\)。
\end{enumerate}

\begin{enumerate}
\def\labelenumi{\alph{enumi})}
\item
  若不相交的集合 \(\mathbf{A}_1\) 与 \(\mathbf{A}_2\) 都是收敛集，则 \(\mathbf{A}_1 \cup \mathbf{A}_2\) 也是收敛集。
\item
  设 A 是一个开集或闭集。为使 A 是收敛集，必要且充分的是 \(\lim_{\lambda,\mathfrak{F}} \lambda^{\bullet}(A) = \mu^{\bullet}(A)\)。若 A 是收敛集，则 \(T - A\) 也是收敛集。
\item
  设 A 是 T 中一个开或闭的收敛集，B 是一个收敛集且 T - B 也是收敛集。则 A ∪ B 与 A ∩ B 都是收敛集，它们的余集亦然。
\item
  由开收敛集生成的集环由收敛集组成。
\item
  设 \(\mathbf{A}\) 是 \(\mathbf{T}\) 的一个边界为 \(\mu\) -可忽略的子集。则 \(\mathbf{A}\) 是一个收敛集。
\item
  假设 T 是局部紧的或波兰的。证明 T 的每个紧子集 K 都包含在一个紧收敛集中（当 T 局部紧时，利用 Exer. 1 b)。若 T 是波兰空间，则利用同一习题构造 K 的一列有限覆盖 \(U_{p}\)（由 T 的开子集组成，边界 \(\mu\) -可忽略且直径 \(<2^{-p}\)）；证明 \(L=\bigcap_{p}\bigcup_{U\in\mathfrak{U}_{p}}\overline{U}\) 是紧的，并由 b) 断定它是一个收敛集）。
\end{enumerate}


\begin{enumerate}
\def\labelenumi{\alph{enumi})}
\setcounter{enumi}{6}
\tightlist
\item
  把 f) 的结果推广到完备度量空间上收敛测度序列的情形。
\end{enumerate}

\begin{enumerate}
\def\labelenumi{\arabic{enumi})}
\setcounter{enumi}{2}
\tightlist
\item
  设 \(\mathbf{T}\) 是一个波兰空间，\((\mu_n)\) 是 \(\mathbf{T}\) 上一列紧收敛到测度 \(\mu\) 的有界正测度。用 \(\mathfrak{C}\) 表示具有下列性质的子集 \(\mathbf{A}\)（\(\mathbf{T}\) 的）所成的集合：对每个 \(\varepsilon > 0\)，存在紧子集 \(\mathbf{K}\)（\(\mathbf{A}\) 的），使得 \(\sup_{n} \mu_n(\mathbf{A} - \mathbf{K}) < \varepsilon\)。用 \(\mathfrak{D}\) 表示子集 \(\mathbf{A}\)（\(\mathbf{T}\) 的，与其余集一起都属于 \(\mathfrak{C}\)）所成的集合。
\end{enumerate}

\begin{enumerate}
\def\labelenumi{\alph{enumi})}
\item
  若集合 A 与 A' 属于 C，则 \(A \cup A'\) 与 \(A \cap A'\) 也属于 C；由此推出 D 是 T 的子集所成的一个集环。
\item
  T 的每个边界为 \(\mu\) -可忽略的子集 A 都属于 \(\mathfrak{D}\)（对 A 的内部应用 Prokhorov 定理，并利用 Exer. 2 e)）。
\item
  每个 \(\mathbf{A} \in \mathfrak{D}\) 都是序列 \((\mu_n)\) 的收敛集（cf.~Exer. 2）；反之，每个开或闭的收敛集都属于 \(\mathfrak{D}\)。
\item
  设 \(\mathbf{A} \in \mathfrak{D}\)。证明存在一列两两不相交的紧集 \(\mathbf{K}_p\) 与子集 \(\mathbf{N}\)（\(\mathbf{T}\) 的），它们具有下列性质：\(\alpha\) ) 每个 \(\mathbf{K}_p\) 都是序列 \((\mu_n); \beta)\) 的收敛集；\(\mathbf{A} \subset \mathbf{N} \cup \bigcup_{p} \mathbf{K}_p\) 且 \(\bigcup_{p} \mathbf{K}_p \subset \mathbf{A} \cup \mathbf{N}; \gamma)\) 对每个 \(\varepsilon > 0\)，存在开邻域 \(\mathbf{U}\)（\(\mathbf{N}\) 的），使得 \(\sup_{n}\mu_n(\mathbf{U}) < \varepsilon\)（对 \(\mathbf{T}\) 的一个适当的子空间（它是一列开集的交）应用 Exer. 2 f) 与 Prokhorov 定理）。
\end{enumerate}

\begin{enumerate}
\def\labelenumi{\arabic{enumi})}
\setcounter{enumi}{3}
\item
  设 \(\mathbf{T}\) 是一个完全正则空间，\(t\) 是 \(\mathbf{T}\) 的一个点，\(\mathfrak{U}\) 是 \(\mathbf{T}\) 的一族 Borel 子集，它生成 \(t\) 的邻域滤子。设 \(\mathbf{I}\) 是带有一个滤子 \(\mathfrak{F}\) 的集合，\((\mu_i)_{i \in \mathbb{I}}\) 是 \(\mathbf{T}\) 上一族总质量为 1 的有界正测度。为使 \(\lim_{i, \mathfrak{F}} \mu_i = \varepsilon_t\)，必要且充分的是 \(\lim_{i, \mathfrak{F}} \mu_i(\mathrm{U}) = 1\)（对每个 \(\mathrm{U} \in \mathfrak{U}\)）。
\item
  设 E 是一个实 Hilbert 空间，具有一个标准正交基 \((x_{n})_{n\in \mathbf{N}}\)。设 T 是赋以弱化拓扑的空间 E，并设 \((a_{n})_{n\in \mathbf{N}}\) 是一列实数，使得 \(0 < a_{n} < 1\)（对所有 \(n\)），且 \(\lim_{n\to \infty}n^2\log a_n = 0\)。对每个 \(n\in \mathbf{N}\)，置 \(\mu_{n} = \sum_{p\in \mathbf{N}}a_{n}^{p}(1 - a_{n})\cdot \varepsilon_{n\cdot x_{p}}\)。证明 \(\mu_{n}\) 对每个 \(n \in \mathbf{N}\) 都是 T 上总质量为 1 的有界正测度，\(\lim_{n\to\infty}\mu_{n}=\varepsilon_{0}\)（紧收敛），并且 \(\lim_{n\to\infty}\mu_{n}(\mathrm{K})=0\)（对 T 的每个紧子集 \(\mathrm{K}\)）（对 0 的邻域组成的集合 \(\mathfrak{U}\)（形如 \(\{x \mid (a \mid x) \leqslant 1\}\)，其中 a 取遍 E）应用上一习题的判别法）。特别地，序列 \((\mu_{n})_{n \in \mathbf{N}}\) 的元素所成的集合是 \(\mathcal{M}_{+}^{b}(\mathbf{T})\) 中一个相对紧子集，但它不满足 Prokhorov 条件。
\end{enumerate}


\begin{enumerate}
\def\labelenumi{\arabic{enumi})}
\setcounter{enumi}{5}
\tightlist
\item
  设 \(\mathbf{T}\) 为一个拓扑空间，\(\mathbf{H}\) 为 \(\mathbf{T}\) 上正测度组成的一个集合；假设下列条件成立：
\end{enumerate}

\begin{enumerate}
\def\labelenumi{(\Alph{enumi})}
\setcounter{enumi}{21}
\tightlist
\item
  \(\mathbf{T}\) 的每个点都有一个开邻域 \(\mathbf{W}\)，使得 \(\sup_{\mu \in \mathbf{H}} \mu^{\bullet}(\mathbf{W})\) 有限。
\end{enumerate}

最后，设 \(\mathfrak{U}\) 为 H 上的一个超滤子。

\begin{enumerate}
\def\labelenumi{\alph{enumi})}
\item
  设 \(\mathbf{K}\) 是 \(\mathbf{T}\) 的一个紧子集。证明诱导测度 \(\mu_{\mathbf{K}} (\mu \in \mathbf{H})\) 关于 \(\mathfrak{U}\) 淡收敛到一个测度 \(\pi^{\mathbf{K}}\)（在 \(\mathbf{K}\) 上）。
\item
  设 \(\mathbf{K}\) 与 \(\mathbf{L}\) 是 \(\mathbf{T}\) 的两个紧子集，且 \(\mathbf{K} \subset \mathbf{L}\)；证明 \((\pi^{\mathbf{L}})_{\mathbf{K}} \geqslant \pi^{\mathbf{K}}\)。
\item
  对 T 的每个紧子集 K，用 \(i^{\mathbf{K}}\) 表示 K 到 T 中的典范单射。证明测度族 \(i^{\mathbf{K}}(\pi^{\mathbf{K}})\) 有一个上确界 \(\pi\)（在 \(\mathcal{M}_{+}(\mathrm{T})\) 中）。
\item
  设 \(f\) 是 \(\mathbf{T}\) 上的一个下半连续正函数，\(g\) 是 \(\mathbf{T}\) 上具有紧支撑的一个上半连续正函数。证明 \(\pi^{\bullet}(f) \leqslant \lim_{\mu, \mathfrak{U}} \mu^{\bullet}(f)\) 与 \(\pi^{\bullet}(g) \geqslant \lim_{\mu, \mathfrak{U}} \mu^{\bullet}(g)\)。
\end{enumerate}

¶ 7) 设 H 是拓扑空间 T 上有界正测度组成的一个集合。假设 \(\sup_{\mu\in H}\mu^{\bullet}(T)\) 有限，并且对每个 \(\varepsilon>0\)，存在 T 的一个紧子集 K

使得 \(\sup_{\mu \in \mathbf{H}}\mu^{\bullet}(\mathrm{T} - \mathrm{K}) < \varepsilon\)。证明上一习题的条件 (V) 是

满足的。设 \(\mathfrak{U}\) 为 H 上的一个超滤子。测度 \(\pi\) 按上一习题中的方式定义。

\begin{enumerate}
\def\labelenumi{\alph{enumi})}
\item
  证明 \(\pi^{\bullet}(g) \geqslant \lim_{\mu, \mathfrak{U}} \mu^{\bullet}(g)\) 对 T 上的每个上半连续有界正函数 \(g\) 成立。由此推出，\(\pi^{\bullet}(f) = \lim_{\mu, \mathfrak{U}} \mu^{\bullet}(f)\) 对 T 上每个有界函数 \(f\)（其不连续点集为 \(\pi\) -可忽略）成立。
\item
  当 T 完全正则时，由 a) 推出 H 对紧拓扑是相对紧的（这为 No.~5 的 Th. 1 在正测度的情形提供了一个新的证明）。
\end{enumerate}

¶ 8) 设 T 是一个完备可分度量空间，d 是它的距离。对 T 的每个闭子集 F 与每个实数 a > 0，用 \(F^{a}\) 表示 \(x \in T\) 中使 \(d(x, F) < a\) 的 x 所成的集合。给定 T 上两个有界正测度 \(\lambda\) 与 \(\mu\)，用 \(D(\lambda, \mu)\) 表示满足下列不等式的实数 a > 0 所成集合的下确界：对 T 的每个闭子集 F，\(\lambda^{\bullet}(F) \leqslant \mu^{\bullet}(F^{a}) + a\)，\(\mu^{\bullet}(F) \leqslant \lambda^{\bullet}(F^{a}) + a\)。

\begin{enumerate}
\def\labelenumi{\alph{enumi})}
\item
  证明 D 是 \(\mathcal{M}_+^b\) 上的一个距离，且 \(\mathrm{D}(\varepsilon_x,\varepsilon_y) = d(x,y)\) 对 T 的任意两个点 \(x\) 与 \(y\)（满足 \(d(x,y) < 1\)）成立。
\item
  设 \(f\) 是 \(T\) 到 \(T\) 中的一个 Borel 映射（cf.~TG, IX, §6, No.~3 或 GT, IX, §6, Exer. 16），\(\lambda\) 是 \(T\) 上的一个有界正测度。证明存在 T 上的一个有界正测度 \(\mu\)，使得对 T 的每个 Borel 子集 A 有 \(\mu^{\bullet}(\mathrm{A}) = \lambda^{\bullet}\big(f^{-1}(\mathrm{A})\big)\)。设 \(a > 0\)，并用 H 表示 \(x \in \mathrm{T}\) 中使 \(d(x, f(x)) \geqslant a\) 的 x 所成的集合；证明 H 是 T 中的 Borel 集，且 \(\mathrm{D}(\lambda, \mu) \leqslant \sup (a, \lambda(\mathrm{H}))\)（对 T 的每个闭子集 F，有 \(\mathrm{F} \cap (\mathrm{T} - \mathrm{H}) \subset f(\mathrm{F}^a)\) 且 \(\lambda^{\bullet}(\mathrm{F}) \leqslant \lambda^{\bullet}(\mathrm{F} \cap (\mathrm{T} - \mathrm{H})) + \lambda^{\bullet}(\mathrm{H})\)）。
\item
  设 \(\mu\) 与 \(\mu_n\)（对 \(n \geqslant 1\)）是 \(\mathbf{T}\) 上的有界正测度，使得 \(\lim_{n \to \infty} \mathrm{D}(\mu_n, \mu) = 0\)。证明序列 \((\mu_n)\) 紧收敛到 \(\mu\)（证明 \(\lim_{n \to \infty} \mu_n^\bullet(\mathrm{T}) = \mu^\bullet(\mathrm{T})\)，且 \(\mu^\bullet(\mathrm{F}) \geqslant \limsup_{n \to \infty} \mu_n^\bullet(\mathrm{F})\)（对每个闭子集 \(\mathbf{F}\)（\(\mathbf{T}\) 的））；再按 §2, No.~6 的 Lemma 3 的方法推出 \(\mu^\bullet(f) \leqslant \liminf_{n \to \infty} \mu_n^\bullet(f)\)（对每个下半连续函数 \(f \geqslant 0\)））。
\item
  反之，证明对每个有界正测度序列 \(\mu_n\)（紧收敛到 \(\mu\)，在 \(\mathcal{M}_+^b\) 中），有 \(\lim_{n \to \infty} D(\mu_n, \mu) = 0\)。可以按如下方式进行：设 \(\varepsilon > 0\)，并设 \(K\) 是一个紧子集（\(T\) 的），使得 \(\mu^\bullet(T - K) < \varepsilon\) 且 \(\sup_{n} \mu_n^\bullet(T - K) < \varepsilon\)；
\end{enumerate}

构造一个有限族 \((\mathrm{B}_{i})_{1\leqslant i\leqslant p}\)（由边界 \(\mu\) -可忽略、直径 \(\leqslant\varepsilon/2\)、两两不相交且并为包含 K 的集合组成）（cf.~Exer. 1）。设 f 是 T 到 T 中的一个映射，它在每个集合 \(B_{1},\ldots,B_{p}\) 与 \(\mathbb{C}(B_{1}\cup\cdots\cup B_{p})\) 上取常值，且 \(f(\mathrm{B}_i) \subset \mathrm{B}_i\)（对 \(1 \leqslant i \leqslant p\)）。用下列条件定义测度 \(\pi_n\) 与 \(\pi\)：对 T 的每个 Borel 子集 A，\(\pi_n^\bullet(\mathrm{A}) = \mu_n^\bullet(\overline{f}(\mathrm{A}))\) 且 \(\pi^\bullet(\mathrm{A}) = \mu^\bullet(f^{-1}(\mathrm{A}))\)；由 b) 推出关系式 \(\mathrm{D}(\pi_n, \mu_n) \leqslant \varepsilon\)，\(\mathrm{D}(\pi, \mu) \leqslant \varepsilon\)，并证明 \(\lim_{n \to \infty} \mathrm{D}(\pi_n, \pi) = 0\)。

\begin{enumerate}
\def\labelenumi{\alph{enumi})}
\setcounter{enumi}{4}
\tightlist
\item
  \(M_{+}^{b}\) 上的距离 D 与紧拓扑相容（注意 \(M_{+}^{b}\) 对紧拓扑是可度量化的）。
\end{enumerate}

¶ 9) 记号与假设同上一习题。设 \((\mu_n)\) 是 \(\mathcal{M}_+^b\) 中对距离 D 的一个 Cauchy 序列；证明序列 \((\mu_n)\) 是收敛的（这给出了空间 \(\mathcal{M}_+^b\) 对紧拓扑是波兰空间这一事实的一个新证明）。可以按如下方式进行：

\begin{enumerate}
\def\labelenumi{\alph{enumi})}
\item
  设 \(\varepsilon > 0\) 与 \(a > 0\) 是两个实数；存在一个有限子集 \(F\)（\(T\) 的），使得 \(\sup_{n} \mu_n^\bullet(T - F^a) \leqslant \varepsilon\)（选取一个整数 \(N \geqslant 1\) 与一个紧子集 \(K\)（\(T\) 的），使得 \(\sup_{n \geqslant N} D(\mu_n, \mu_N) \leqslant \varepsilon/2\) 且 \(\mu_N^\bullet(T - K) \leqslant \varepsilon/2\)；由此推出 \(\sup_{n \geqslant N} |\mu_n^\bullet(T) - \mu^\bullet(T)| \leqslant \varepsilon\) 与 \(\sup_{n \geqslant N} \mu_n^\bullet(T - K^{a/2}) \leqslant \varepsilon\)；最后选取一个有限子集 \(F\)（\(T\) 的），使得 \(K^{a/2} \subset F^a\) 且 \(\sup_{n \geqslant N} \mu_n^\bullet(T - F^a) \leqslant \varepsilon\)）。
\item
  设 \(\varepsilon > 0\)；存在一个紧子集 \(K\)（\(T\) 的），使得 \(\sup_{n} \mu_n^\bullet(T - K) \leqslant \varepsilon\)（对每个整数 \(p \geqslant 1\)，选取一个有限子集 \(F_p\)（\(T\) 的），使得 \(\mu_n^\bullet(T - (F_p)^{2^{-p}}) \leqslant \varepsilon / 2^p\)，并置 \(K = \bigcap_{p\geqslant 1}(\overline{(F_p)^{2^{-p}}})\)。
\item
  由 \(b\) ) 推出，Cauchy 序列 \((\mu_n)\) 有一个收敛的子序列（对距离 D）。
\end{enumerate}

¶ 10) 设 T 是一个完全正则空间；用 S 表示 T 上满足下列性质的实函数 \(f \geqslant 1\) 所成的集合：对每个实数 c，使 \(f(t) \leqslant c\) 的 T 的点 t 所成的集合是紧的。~对每个 \(f \in S\)，用 \(\mathcal{M}_f\) 表示 T 上有界测度 \(\mu\)（满足 \(|\mu|^{\bullet}(f) \leqslant 1\)）所成的集合。

\begin{enumerate}
\def\labelenumi{\alph{enumi})}
\item
  为使 \(\mathcal{M}^{b}(T)\) 的子集 H 满足 Prokhorov 条件，必要且充分的是存在 \(f \in S\)，使得 \(H \subset M_{f}\)（对必要性，取 f 形如 \(c(1 + \sum_{n \geqslant 1} n \cdot f_{n})\)，其中 \(f_{n}\) 是这样的集合 \(U_{n}\) 的特征函数：\(\mathrm{T} - \mathrm{U}_n\) 是紧的，且 \(\sup_{\mu \in \mathbf{H}} |\mu|^{\bullet} (\mathrm{U}_n) \leqslant 2^{-n}\)。
\item
  假设 T 是局部紧的。为使 \(\mathcal{M}^{b}(T)\) 的子集 H 是相对紧的（对紧拓扑），必要且充分的是存在连续函数 \(f \in S\)，使得 \(H \subset M_{f}\)。
\item
  设 C 是 \(\mathcal{M}_{+}^{b}(\mathrm{T})\) 中一个闭凸锥。证明 \(\mathrm{C} \cap \mathcal{M}_{f}\) 对每个 \(f \in \mathcal{S}\) 都是 C 的一个帽（TVS, II, §7, No.~2, Def. 3），并由此推出 C 是它的帽的并。用 E 表示 C 的极端生成元的并。假设 T 是 Souslin 空间。证明对每个 \(\pi \in \mathrm{C}\)，存在 \(\mathcal{M}_{+}^{b}(\mathrm{T})\) 中包含于 E 的一个 Borel 子集 B 与 B 上一个总质量为 1 的正测度 P，使得 \(\pi = \int_{\mathrm{B}} \mu d\mathrm{P}(\mu)\)（应用 Choquet 积分表示定理）。
\end{enumerate}

\begin{enumerate}
\def\labelenumi{\arabic{enumi})}
\setcounter{enumi}{10}
\item
  设 T 是一个完全正则空间。假设对每个紧空间 K 与每个 T 到 K 中的连续映射 f，都给定 K 上的一个正测度 \(\mu_{f,K}\)；假设 \(g(\mu_{f,K}) = \mu_{g \circ f,L}\)（对任意紧空间 K 与 L 以及任意连续映射 \(f : T \to K\) 与 \(g : K \to L\)）；再假设对任意 f 与 K，测度 \(\mu_{f,K}\) 集中在 \(f(T)\) 上。证明存在唯一的 T 上的有界正测度 \(\pi\)，使得对任意 f 与 K 有 \(\mu_{f,K} = f(\pi)\)（利用 T 的 Stone--Čech 紧化的泛性质）。
\item
  设 T 是一个完全正则空间，是一个局部紧空间 L 的子空间。对 T 上的每个测度 \(\mu\)（无论正与否），存在 L 的一个包含 T 的局部紧子空间 \(T'\)，以及一个集中在 T 上的测度 \(\mu'\)（在 \(T'\) 上），它在 T 上诱导出 \(\mu\)。
\end{enumerate}

¶ 13) *设 G 是一个局部紧交换群。用 \(\widehat{G}\) 表示 G 的对偶群，用 \(\alpha\) 表示 \(\widehat{G}\) 上的一个 Haar 测度。对 G 上的每个有界测度 \(\mu\)，用 \(\mathcal{F}\mu\) 表示 \(\mu\) 的 Fourier 变换，它是群 \(\widehat{G}\) 上的一个函数（cf.~Theor. spec., Ch. II, §1, No.~2）。

\begin{enumerate}
\def\labelenumi{\alph{enumi})}
\item
  Fourier 变换 \(\mathcal{F}\) 是从赋以紧拓扑的 \(\mathcal{M}_{+}^{b}(\mathrm{G})\) 到赋以紧收敛拓扑的 \(\mathcal{C}^b (\widehat{\mathrm{G}})\) 中的一个连续映射。（利用 No.~6 的 Prop. 13 的推论。）
\item
  设 \(\mathfrak{F}\) 是 \(\mathcal{M}_+^b (\mathrm{G})\) 上的一个滤子，\(\Phi\) 是 \(\widehat{\mathrm{G}}\) 上的一个有界连续函数；假设 \(\lim_{\lambda ,\mathfrak{F}}(\mathcal{F}\lambda)\cdot \alpha = \Phi \cdot \alpha\)（淡收敛）且 \(\lim_{\lambda ,\mathfrak{F}}(\mathcal{F}\lambda)(0) = \Phi (0)\)。证明滤子 \(\mathfrak{F}\) 紧收敛到一个测度 \(\mu\)，使得 \(\mathcal{F}\mu = \Phi\)（应用 Stone-Weierstrass 定理证明：G 上具有紧支撑的连续函数的 Fourier 变换所成的集合，是 \(\mathcal{C}^0 (\overline{\mathrm{G}})\) 中对一致收敛的稠密线性子空间；接着注意 \(\lim_{\lambda ,\mathfrak{F}}\int (\mathcal{F}\lambda)\cdot u d\alpha = \int \Phi u d\alpha\) 对每个 \(u\in \mathrm{E}\) 成立，且滤子 \(\mathfrak{F}\) 含有有界测度的一个淡紧集合；由此推出滤子 \(\mathfrak{F}\) 至多有一个聚点，然后推出它紧收敛）。
\item
  设 \(\mathfrak{F}\) 是 \(\mathcal{M}_{+}^{b}(\mathrm{G})\) 上具有可数基的一个滤子。为使 \(\mathfrak{F}\) 紧收敛到一个有界正测度 \(\mu\)，必要且充分的是存在一个连续函数 \(\Phi\)（在 \(\widehat{\mathbf{G}}\) 上），使得 \(\mathcal{F}\lambda\) 逐点收敛到 \(\Phi\)（关于 \(\mathfrak{F}\)），此时 \(\mathcal{F}\mu = \Phi\)（\(P\) . Lévy 定理）。*
\end{enumerate}

